
\section{AA-03: Cesium clock splitting and weak-induced E1 amplitude}
\source{Week 13 clock slides; Week 12 atomic-parity-violation slides.}
\subsection*{Example A: the clock's hyperfine states}
For \(^{133}\mathrm{Cs}\) ground state \({}^2S_{1/2}\),
\(J=1/2\) and nuclear spin \(I=7/2\). Therefore
\[
\boxed{F=3,4},\qquad
I(I+1)=63/4,\quad J(J+1)=3/4.
\]
Using the magnetic-dipole hyperfine formula,
\[
\begin{aligned}
E_3&=\frac A2\left(12-\frac{63}{4}-\frac34\right)
=\boxed{-\frac{9A}{4}},\\
E_4&=\frac A2\left(20-\frac{63}{4}-\frac34\right)
=\boxed{\frac{7A}{4}} .
\end{aligned}
\]
The clock frequency is determined by
\[
\boxed{E_4-E_3=4A=h\nu_{\rm Cs}},\qquad
\nu_{\rm Cs}=9\,192\,631\,770\ {\rm Hz}.
\]
Consequently \(A/h=\nu_{\rm Cs}/4
\simeq2.29816\ \mathrm{GHz}\).
The actual clock transition is
\(\ket{F=3,m_F=0}\leftrightarrow\ket{F=4,m_F=0}\):
both \(m_F=0\) states have no linear Zeeman shift
in the weak-field approximation, reducing
magnetic-field sensitivity.
\(\sum_F(2F+1)=7+9=16=(2I+1)(2J+1)\)
checks the state count.

\subsection*{Example B: weak interaction opens a forbidden line}
Suppose \(\ket i\) and \(\ket f\) have the \emph{same} parity.
An electric dipole \(d\) is parity-odd, therefore
\(\langle f|d|i\rangle=0\) if parity is exact.
Introduce a weak parity-violating perturbation \(H_W\)
that mixes the initial state with an opposite-parity
state \(\ket n\):
\[
\ket{\tilde i}\simeq\ket i+
\epsilon\ket n,\qquad
\epsilon=\frac{\langle n|H_W|i\rangle}{E_i-E_n}.
\]
To first order,
\[
\boxed{\langle f|d|\tilde i\rangle
\simeq\epsilon\langle f|d|n\rangle}.
\]
For an \emph{illustrative} ratio
\(|\epsilon|=10^{-8}\) and allowed
\(|\langle f|d|n\rangle|=1\ e a_0\),
the induced amplitude is only
\(10^{-8}ea_0\); a measurement needs a
sensitive interference technique.
In general, corrections to both initial and final
states contribute, with a sum over intermediate states.

\subsection*{Why this connects to EDM experiments}
An intrinsic permanent EDM along angular momentum
would violate parity and time reversal.
An external electric field gives
\(H_{\rm EDM}=-\mathbf d\cdot\mathbf E\).
Field reversal is used to distinguish an EDM-like
signal from many conventional shifts. Atomic
structure and symmetry selection rules are thus
essential for precision fundamental-physics tests.

\checkmethod{Do not confuse a \emph{clock}
transition frequency with the frequency of an
allowed optical E1 line. Weak perturbations can
produce tiny transitions that are exactly forbidden
in the unperturbed parity basis.}
